\section{Требования к пользовательскому интерфейсу}

Основные пользователи веб-интерфейса~--- студенты: они загружают свои материалы (шаблоны ВКР или отчёты по практике) и получают отчёты с результатами формальной проверки (предупреждения, ошибки). Дополнительно доступен интерфейс на основе электронной почты (раздел~3.5).

\subsection{Организация по типу документа}

Интерфейс организован по типу документа: \emph{шаблон ВКР} и \emph{отчёт по практике}. Списки, фильтры и дашборд оперируют типами документов, статусами и датами.

\subsection{Сценарий загрузки}

Пользователь выбирает тип документа, затем загружает один или несколько PDF-файлов (перетаскивание в зону загрузки или выбор через диалог). После загрузки запускается проверка; по завершении отображаются результаты и формируется отчёт.

\subsection{Результаты проверки}

Список документов отображается в виде таблицы (или карточек) с колонками: имя файла, тип документа, статус (Пройдено / Не пройдено / Предупреждение), дата. По клику открывается детальный отчёт по проверке и просмотр PDF с подсветкой мест, связанных с замечаниями. Реализуется экспорт отчёта о проверке. Система автоматически предоставляет отчёт с результатами; ручная проверка и комментарии преподавателя не предусмотрены.

\subsection{Нормативные документы}

В интерфейсе при необходимости размещаются ссылки на нормативные документы ИТМО по оформлению ВКР и отчётных документов по практике.
